\section{Introduction}

Consider a collection of linear polymers with total mass $n$. Each pair
of polymers merges at a constant rate, and each internal link breaks
independently. The resulting Markov chain has an explicit reversible
law, but that law alone does not determine how fast the chain reaches
equilibrium. Aldous~\cite{AldousNotes,AldousProblemPage} asked for its
relaxation scale when the fragmentation rate is comparable to $1/n$.
We determine that scale with its leading constant, obtain a cutoff
theorem, and study extensions with spatial motion and nonreversible
rates.

The argument begins with the cycle representation in Aldous's notes.
Its labelled cut process has more slow modes than the polymer chain.
Rotation invariance removes the extra degree-one modes, and an explicit
comparison puts every remaining harmonic sector above the first
nonconstant count mode. This gives a direct proof of the equality
between the polymer and count gaps announced by Granovsky and
Erlihson~\cite{GranovskyErlihson2009}. A choice of reciprocal test
function then gives a spectral estimate with an absolute error of
order $1/n$.

\subsection{The chain and the main theorem}

Fix $n\ge2$ and $\beta>0$. The state space is
\[
 \mathcal S_n=\left\{m=(m_1,\ldots,m_n)\in\mathbb Z_{\ge0}^n:
                     \sum_{j=1}^n j m_j=n\right\}.
\]
Here $m_j$ is the number of clusters containing $j$ units of mass. Each cluster is viewed as a linear polymer with $j-1$ internal links. Every unordered pair of distinct cluster instances merges at rate $2/n$; the two instances may have the same size. Every internal link breaks at rate $\beta$.

Let $e_j$ be the $j$th coordinate vector. The nonzero off-diagonal transition rates of the generator $L_M$ are
\begin{align}
 q_M(m,m-e_j+e_i+e_{j-i})
  &=\begin{cases}
        2\beta m_j,&1\le i<j/2,\\
        \beta m_j,&i=j/2,
     \end{cases}\label{eq:fragmentation-rates}\\
 q_M(m,m+e_{i+j}-e_i-e_j)
  &=\begin{cases}
        2m_i m_j/n,&i<j,\quad i+j\le n,\\
        m_i(m_i-1)/n,&i=j,\quad 2i\le n.
     \end{cases}\label{eq:coagulation-rates}
\end{align}
Rates with an inadmissible target are zero. The diagonal entries are the negatives of the row sums. The factors of two in~\eqref{eq:fragmentation-rates} count the two distinct links giving the same unordered split when the two fragments have different sizes.

Write
\[
 K(m)=\sum_{j=1}^n m_j
\]
for the number of clusters. The chain is irreducible and reversible with respect to
\begin{equation}\label{eq:polymer-stationary-intro}
 \pi_M(m)=\frac1{Z_M}\,
                 \frac{(\beta n)^{K(m)}}{\prod_{j=1}^n m_j!}.
\end{equation}
The normalizing constant $Z_M$ is the sum of the displayed weights over $\mathcal S_n$. We rederive reversibility below, together with the comparison used in the proof.

For real functions on $\mathcal S_n$, let
\begin{equation}\label{eq:dirichlet-intro}
 \mathcal E_M(f,f)
 =-\langle f,L_Mf\rangle_{\pi_M}
 =\frac12\sum_{m,m'\in\mathcal S_n}
       \pi_M(m)q_M(m,m')\bigl(f(m')-f(m)\bigr)^2.
\end{equation}
The spectral gap and relaxation time are
\begin{equation}\label{eq:gap-intro}
 \gamma_M(n,\beta)=
  \inf_{\operatorname{Var}_{\pi_M}f>0}
       \frac{\mathcal E_M(f,f)}{\operatorname{Var}_{\pi_M}f},
 \qquad
 \tau_{\mathrm{rel}}(n,\beta)=\gamma_M(n,\beta)^{-1}.
\end{equation}
We usually abbreviate $\gamma_M(n,\beta)$ to $\gamma_M$. All constants implicit in $\asymp$ below are independent of $n$ and $\beta$.

\begin{theorem}\label{thm:main}
For every $n\ge2$ and $\beta>0$, set
\[
 s(n,\beta)=\beta+\sqrt\beta+\frac1n.
\]
Then
\begin{equation}\label{eq:main-gap}
 \frac{s(n,\beta)}6\le\gamma_M(n,\beta)\le4s(n,\beta).
\end{equation}
Consequently,
\begin{equation}\label{eq:main-relaxation}
 \frac1{4s(n,\beta)}\le\tau_{\mathrm{rel}}(n,\beta)
                         \le\frac6{s(n,\beta)}.
\end{equation}
The proof also gives the estimates
\begin{equation}\label{eq:sharper-elementary}
 \beta+\frac2n\le\gamma_M(n,\beta)
       \le\beta+4\max\!\left\{\sqrt\beta,\frac1n\right\}.
\end{equation}
\end{theorem}

Theorem~\ref{thm:main} remains useful at the transition
$\beta\asymp n^{-2}$, where the additive error in the sharper estimate
below is comparable to the gap. The one-state case $n=1$ has no
positive eigenvalue and is excluded throughout.

\begin{samepage}
\begin{corollary}\label{cor:regimes}
Let $\beta_n>0$ and let $n\to\infty$.
\begin{enumerate}
 \item If $\beta_n\le Cn^{-2}$ for a fixed $C<\infty$, then
       $\tau_{\mathrm{rel}}(n,\beta_n)\asymp n$.
 \item If $n^{-2}\ll\beta_n\ll1$, then
       $\tau_{\mathrm{rel}}(n,\beta_n)\asymp\beta_n^{-1/2}$.
 \item If $\beta_n\ge c>0$, then
       $\tau_{\mathrm{rel}}(n,\beta_n)\asymp\beta_n^{-1}$.
\end{enumerate}
The comparison constants in the first and third assertions may depend on $C$ and $c$, respectively. In particular, $\beta_n\asymp n^{-1}$ implies
$\tau_{\mathrm{rel}}(n,\beta_n)\asymp\sqrt n$.
\end{corollary}
\end{samepage}

\begin{proof}
In each range, compare the three terms in $s(n,\beta_n)$ and apply~\eqref{eq:main-relaxation}. If $\beta_n\le Cn^{-2}$, the square-root term is at most $\sqrt C/n$. In the second range it dominates both other terms. If $\beta_n\ge c$, then $\sqrt{\beta_n}\le c^{-1/2}\beta_n$ and $n^{-1}\le(2c)^{-1}\beta_n$.
\end{proof}

For fixed $\beta$, Aldous's estimate is
\begin{equation}\label{eq:aldous-bound}
 (1-o(1))(\beta^2+4\beta)^{-1/2}
 \le\tau_{\mathrm{rel}}(n,\beta)\le\beta^{-1}
 \qquad(n\to\infty).
\end{equation}
His notes also discuss $\beta_n\sim b/n$, where the corresponding bounds
have different powers of $n$~\cite{AldousNotes}. The uniform theorem
closes that order gap. The sharper results identify the constant in
both regimes.

\subsection{Sharp conclusions and extensions}

Let $\gamma_Q$ be the gap of the cluster-count chain. We prove
\[
 \gamma_M=\gamma_Q,\qquad
 \sqrt{\beta^2+4\beta}-\frac5n
 \le\gamma_M\le
 \sqrt{\beta^2+4\beta}+\frac4n
\]
in Theorems~\ref{thm:exact-gap} and~\ref{thm:sharp-gap}.
It follows that
\[
 \tau_{\mathrm{rel}}(n,b/n)\sim\frac{\sqrt n}{2\sqrt b},
 \qquad
 \tau_{\mathrm{rel}}(n,\beta)\longrightarrow
       (\beta^2+4\beta)^{-1/2}
 \quad\text{for fixed }\beta>0.
\]
At the remaining transition scale we construct a positive, continuous
function $\rho$ such that
\[
 \left|n\gamma_M(n,b/n^2)-\rho(b)\right|\le b/n.
\]
The function is the lowest eigenvalue of an explicit infinite
birth--death operator. Finite tridiagonal matrices approximate it with
a proved truncation error; see Section~\ref{sec:critical-scaling}.

For $\beta=b/n$ with $b>0$ fixed, the count process has worst-case
total-variation cutoff at
\[
 t_n=\frac{\sqrt n\log n}{8\sqrt b},
\]
with a window of order at most $\sqrt n$
(Theorem~\ref{thm:count-cutoff}). The same statement holds for the
polymer process from either extreme initial configuration: one polymer
of size $n$, or $n$ monomers. This transfer is exact at every time,
because the density from either state remains a function of the
cluster count. Separation cutoff from these states occurs at $2t_n$.
For arbitrary initial partitions,
Theorem~\ref{thm:full-polymer-tv-upper-new} gives
\[
 t^{\mathrm{polymer}}_{\mathrm{mix},n}(\varepsilon)
 \le C_b\sqrt n\bigl[(\log n)^2+\log(1/\varepsilon)\bigr].
\]
Thus the worst-case polymer mixing time is determined within one
logarithmic factor. This upper bound uses positivity of the actual
polymer density and the degree-dependent sector energies.

The cycle representation also leads to a spatial process on cut sets.
Add symmetric exclusion along a connected weighted graph $G$ with
diffusion coefficient $D$. Theorem~\ref{thm:spatial-exact-gap} gives
\[
 \gamma_{\mathrm{sp}}
   =\min\{\gamma_Q,\eta_1+D\lambda_G\},
\]
where $\lambda_G$ is the one-particle graph gap and $\eta_1$ is the
lowest eigenvalue of the degree-one radial reaction operator.
Writing $d=\sqrt{\beta^2+4\beta}$ and $c_1=(\beta+d)/2$, we obtain
\[
 \min\{d,c_1+D\lambda_G\}-\frac5n
 \le\gamma_{\mathrm{sp}}\le
 \min\{d,c_1+D\lambda_G\}+\frac4n.
\]
An additional critical scaling function describes $\eta_1$ when
$\beta\asymp n^{-2}$. On a cycle the dynamics moves interval
boundaries and merges or splits neighboring intervals. On a general
graph it is an exclusion process with creation and deletion rates
depending on the total occupancy.

Finally, bounded reversible changes of the reaction conductances
give uniform estimates for inhomogeneous rates. Conservative cut
motion generated by directed permutations has an exact formula for
its asymptotic spectral rate and a separate formula for its
$L^2$ contraction rate. A stationary current on four existing
polymer configurations strictly increases the asymptotic spectral
rate of the original partition-valued chain. These extensions are
proved in Sections~\ref{sec:spatial-diffusion}--\ref{sec:kernel-comparison}.

\subsection{Relation to earlier work}

Reversible clustering processes and their product-form equilibrium
laws are classical; see Kelly~\cite[Chapter~8]{Kelly1979} and
Whittle~\cite{Whittle1986}. Durrett, Granovsky, and
Gueron~\cite{DurrettGranovskyGueron1999} characterize reversibility
for a broad class of coagulation--fragmentation processes and study
the resulting equilibrium distributions. Berestycki and
Pitman~\cite{BerestyckiPitman2007} relate Gibbs laws on partitions
to recursive fragmentation and its coalescent time reversal.
The stationary law used here belongs to this established setting.

The solvable family of Granovsky and
Erlihson~\cite{GranovskyErlihson2009} contains the present model.
In their notation, set
\[
 N=n,\qquad a=0,\qquad b=2/n,\qquad\phi(1,1)=\beta.
\]
Their Proposition~1 gives the autonomous count process, and their
fragmentation formulas give the same unequal- and equal-fragment
rates as~\eqref{eq:fragmentation-rates}. Their Section~2.3 announces
equality of the full-process and count gaps and gives the bounds
\[
 \beta+2/n\le\gamma_Q\le\beta+2-2/n.
\]
The displayed transition formula used there concerns initial laws
with Gibbs conditional distributions at fixed count. Our proof of
the gap equality instead controls every orthogonal sector of every
polymer observable. We use this independent comparison to obtain
the uniform error estimate, the critical limit, and the graph
extensions. The gap equality itself is attributed to that earlier
work.

Hoze and Holcman~\cite{HozeHolcman2018} analyze the same constant-kernel
model after a time change: take their fragmentation parameter to be
$\beta n/2$ and multiply their generator by $2/n$. They obtain
equilibrium cluster statistics and quantities for tagged clusters.
D'Orsogna, Lei, and Chou~\cite{DOrsognaLeiChou2015} study first
assembly and equilibration in finite stochastic clustering systems,
including restrictions on the largest allowed cluster. The
all-observable spectral estimate and the mixing statements here
concern the unrestricted partition-valued chain with the stated
normalization.

Our one-dimensional arguments use positive-vector estimates for
birth--death chains, developed in work including
Zeifman~\cite{Zeifman1991} and Granovsky and
Zeifman~\cite{GranovskyZeifman2000,GranovskyZeifman2005}.
The derivative identity is a finite unweighted instance of the
intertwining relations of Chafa\"i and Joulin~\cite{ChafaiJoulin2013}.
Filmus~\cite{Filmus2016} gives the orthogonal harmonic structure on
Boolean slices. We prove the finite identities used in the argument
so that the cut-to-polymer comparison and its boundary conventions
are explicit.

The cutoff step uses the birth--death hitting-time theorems of Chen
and Saloff-Coste~\cite{CSC2015}. The calculation here supplies the
endpoint means and variances for the moving equilibrium scale
$K\asymp\sqrt n$. More general concentration and cutoff results for
density-dependent chains appear in Barbour, Brightwell, and
Luczak~\cite{BBL2022,BBL2026}. A direct application of a
fixed-density theorem would require checking both the
$n$-dependent rescaled rates and the all-monomer starting state,
whose rescaled density diverges. The passage-time proof handles
these features directly.

Spatial coagulation--fragmentation has a separate literature.
Bernardin and Toninelli~\cite{BernardinToninelli2012} study a
one-dimensional interval process with a dynamical phase transition
and mixing estimates. Ioffe and T\'oth~\cite{IoffeToth2020} study
split-and-merge limits for stationary random stirring on a lattice
torus. Our graph model retains exchangeable
reaction rates and adds exclusion motion. It is related to graph
gap reductions for exclusion with reservoirs, including
Salez~\cite{Salez2023,Salez2024}, and to the one- and two-particle
identities of Kim and Sau~\cite{KimSau2026}. The reservoir models
have fixed sitewise creation and annihilation rates, whereas our
deletion rate is $(|S|-1)/n$ and the reversible law is not a product
law. Here the higher reaction sectors already dominate the count
gap, leaving only the one-particle graph calculation in the exact
formula. The proof therefore does not require the general
interchange-process gap theorem of Caputo, Liggett, and
Richthammer~\cite{CaputoLiggettRichthammer2010}.

There are also established growing-kernel and nonreversible
comparisons. Jiang~\cite{Jiang2015} diagonalizes the Metropolis
transposition chain with Ewens equilibrium. We record its spectral
gap in the present time units as a known baseline for a different
multiplicative coagulation rule. Antisymmetric perturbations that
preserve equilibrium can accelerate convergence; Hwang, Hwang-Ma,
and Sheu~\cite{HwangHwangMaSheu2005} give a spectral treatment for
diffusions. Our four-state circulation verifies strict acceleration
on the actual reaction graph of the linear-polymer chain.

Deterministic coagulation--fragmentation equations require different
arguments. Aizenman and Bak~\cite{AizenmanBak1979} prove convergence
through decay of free energy. Recent work of Bhattacharya and
Throm~\cite{BhattacharyaThrom2026} obtains convergence without
detailed balance for perturbations of constant continuous kernels.
The finite graph and kernel results below concern their explicitly
defined Markov generators. They do not supply a convergence theorem
for general spatial coagulation--fragmentation equations.


\subsection{Structure of the proof}

There are two comparisons. First, the cluster count is a birth--death chain $Q$ on $\{1,\ldots,n\}$ with rates
\begin{equation}\label{eq:count-rates-intro}
 \lambda_k=\beta(n-k),\qquad \mu_k=\frac{k(k-1)}n.
\end{equation}
Every count observable is a polymer observable, so the variational principle gives $\gamma_M\le\gamma_Q$. Discrete differentiation turns the positive spectrum of $-Q$ into the spectrum of a killed birth--death operator. A positive function of the form $[(k+a)(k+a+1)]^{-1}$, with $a=\max\{1,n\sqrt\beta\}$, gives the upper bound in~\eqref{eq:sharper-elementary} directly.

For the lower bound, place $n$ bonds in cyclic order and retain a nonempty set of cuts. Add each absent cut at rate $\beta$ and remove each present cut at rate $(k-1)/n$ when there are $k$ cuts. The unordered cyclic segment lengths define a polymer configuration. At stationarity this map preserves the variance and Dirichlet form of every lifted polymer observable. This is the precise comparison needed; the projected paths need not form a Markov chain from every deterministic cut configuration.

The cut-set generator commutes with all permutations of the bonds.
Its harmonic decomposition reduces it to killed birth--death
operators indexed by a degree $\ell$. Degree zero is the count chain.
Every polymer observable is invariant under cyclic rotation, which
annihilates its degree-one component. The remaining sectors first
give the uniform estimate and then, through a comparison with the
positive count-derivative eigenfunction, the exact gap reduction.
Reciprocal functions chosen by algebraic cancellation yield the
sharp constants. Hitting-time means and variances give the cutoff
theorem. For graph diffusion the harmonic sectors are tensor sums;
all degrees above one already lie above the count gap before
diffusion is added. A separate coupling proof of the square-root
estimate is given in the appendix.
